% ==============================================================================
%  Chapter 2 -- Lattice gauge theory
%  Purpose: self-contained LGT primer, ~15-20 pages. Everything later chapters
%  rely on must be defined HERE (links, plaquettes, Wilson action, gauge
%  transformations, transfer matrix, correlators, anisotropy, topology).
%  Sources to draw from: notes/papers_review.md sec. 0 (lattice primer) and
%  glueball_report/glueball_spectroscopy.tex secs. "From the path integral
%  to the lattice" and "Masses from Euclidean correlators".
% ==============================================================================

\section{The Euclidean path integral}
\label{sec:lgt-path-integral}

\TODO{Wick rotation, Euclidean action, path integral as a statistical
partition function $Z = \int \mathcal{D}U\, e^{-S[U]}$; expectation values
of observables. This is what licenses Monte Carlo (Ch.~3) and the
transfer-matrix reading of correlators (\cref{sec:lgt-transfer}).}

\section{Gauge fields on the lattice}
\label{sec:lgt-links}

\subsection{Links, plaquettes, and the Wilson action}
\TODO{Define link variables $U_\mu(x) \in G$ as parallel transporters,
gauge transformations
$U_\mu(x) \to \Omega(x)\, U_\mu(x)\, \Omega^{\dg}(x+\hat\mu)$, the
plaquette
$P_{\mu\nu}(x) = U_\mu(x) U_\nu(x+\hat\mu) U^{\dg}_\mu(x+\hat\nu) U^{\dg}_\nu(x)$
and the Wilson action
$S = \beta \sum_p \left(1 - \tfrac{1}{N_c}\Re\Tr P\right)$ \cite{wilson1974}.
Naive continuum limit. Fix the conventions used by the code (they match
the repo's CLAUDE.md): direction-first tensor layout, periodic boundary
conditions, $\mu < \nu$ plaquette ordering.}

\subsection{Gauge groups used in this thesis}
\TODO{$\mathbb{Z}_2$ as the debug-friendly testbed (why: exact results,
bit-exact float64 invariance tests, cheap), $\SU{2}$ as the first
non-abelian production group, the $\U{1}/\SU{3}$ generalisation path.
Note the design rule that makes the code group-generic: color axes always
present ($n_c = 1$ for $\mathbb{Z}_2$), every product an explicit matrix
product, every inverse an explicit dagger.}

\subsection{Wilson loops and observables}
\TODO{Rectangular $R \times T$ Wilson loops, the static potential /
confinement reading, and the naive (plaquette-based, clover-free)
topological charge density $q(x)$ with
$F_{\mu\nu} = (P - P^{\dg})/2i$ -- these are the regression targets of
\cref{ch:regression}.}

\section{Masses from Euclidean correlators}
\label{sec:lgt-transfer}

\subsection{The transfer matrix and spectral decomposition}
\TODO{Transfer matrix $\hat T = e^{-\at \hat H}$, spectral decomposition of
the two-point function
$C(\Delta) = \sum_n |\langle 0|\hat O|n\rangle|^2 e^{-m_n \Delta}$
(periodic: cosh). Key inequality for Ch.~7: for ANY operator on a single
timeslice, $C(\Delta)/C(0) \le e^{-m_0 \Delta}$ -- the variational bound
the Rayleigh loss saturates. Stress that the bound REQUIRES the operator
to live on one timeslice (this becomes the per-timeslice network
constraint of \cref{sec:glue-constraint}).}

\subsection{Zero momentum projection, vacuum subtraction, effective mass}
\TODO{$\Obar(t) = \sum_{\vec x} O(t, \vec x)$; the $0^{++}$ channel has a
vacuum expectation value, so the connected correlator needs explicit
subtraction. Define $\meff(\Delta) = \log[C(\Delta)/C(\Delta+1)]$ and the
cosh-corrected version; plateau reading.}

\subsection{Glueballs and why they are hard}
\TODO{What a glueball is; the signal-to-noise problem (signal
$\sim e^{-m\Delta}$ against a $\Delta$-independent vacuum noise floor);
why heavy masses in lattice units are unresolvable. Sets up smearing,
GEVP, and anisotropy as the classical rescue toolkit
(\cref{sec:glue-classical}).}

\section{Anisotropic lattices}
\label{sec:lgt-anisotropy}

\TODO{Fine temporal spacing $\at = \as/\xi$ buys more correlator points
per physical decay length -- the standard trick for heavy states
\cite{morningstar1999}. Tree-level anisotropic Wilson action:
$\beta_t = \beta\xi$ on temporal plaquettes, $\beta_s = \beta/\xi$ on
spatial ones. Bare vs.\ renormalized anisotropy (the code measures
$\xi_R$ via a Creutz-ratio ratio in \texttt{validate\_anisotropy.py});
be explicit that this thesis uses the BARE tree-level $\xi$ and does not
tune it -- a declared limitation (\cref{sec:concl-future}).}
